\documentclass[11pt,a4paper]{article}

\usepackage[english, bulgarian]{babel}
\usepackage[T2A]{fontenc}
\usepackage{amssymb}
\usepackage{amsmath}
\usepackage{fancyhdr}

\ifdefined\showmargins
\PassOptionsToPackage{showframe}{geometry}
\fi
\usepackage[top=2cm, bottom=2cm, left=2cm, right=2cm]{geometry}

\pagestyle{fancy}
\fancyhf{}
\fancyhead[C]{СДРУЖЕНИЕ С НЕСТОПАНСКА ЦЕЛ „ПЕШАЧКО”}
\fancyhead[R]{\thepage}
\setlength{\headheight}{15pt}% to suppress fancyhdr warnings

\renewcommand{\labelenumi}{(\arabic{enumi})}

\newcounter{article}
\newcommand{\article}[1]{
  \refstepcounter{article}% increment counter
  \noindent\textbf{Чл. \arabic{article}.}
  \, #1
}

\newcommand{\heading}[1]{
  \begin{center}
    \vspace{0.2cm}
    {\bf\large \MakeUppercase{#1}}
  \end{center}
  \addcontentsline{toc}{subsection}{#1}
}

\renewcommand{\thesection}{\Roman{section}}

\begin{document}

\begin{center}
  {\bf\Huge УСТАВ} \\
  на \\
  сдружение с нестопанска цел „ПЕШАЧКО”
\end{center}

\tableofcontents

\vspace{0.7cm}

\section{Наименование, цели, средства, предмет}

\heading{Статут}
\article{}
\begin{enumerate}
\item „ПЕШАЧКО” е сдружение (наричано по-долу „СДРУЖЕНИЕТО”) с нестопанска цел за
  осъществяване на дейност в обществена полза по смисъла на Закона за юридическите лица
  с нестопанска цел (ЗЮЛНЦ).
\item СДРУЖЕНИЕТО е юридическо лице отделно от членовете си и отговаря за задълженията
  си със своето имущество. Членовете на СДРУЖЕНИЕТО не отговарят лично за неговите
  задължения.
\end{enumerate}

\heading{Наименование}
\article{}
\begin{enumerate}
\item Наименованието на СДРУЖЕНИЕТО е: СДРУЖЕНИЕ С НЕСТОПАНСКА ЦЕЛ „ПЕШАЧКО”, което ще
  се изписва на латиница „PESHACHKO”.
\item Всяко писмено изявление от името на СДРУЖЕНИЕТО трябва да съдържа данните описани
  в {\bf Чл. 9} от ЗЮЛНЦ.
\end{enumerate}

\heading{Седалище и адрес}
\article{Седалището на СДРУЖЕНИЕТО е в Република България, гр. София, а адресът на
  управление -- гр. София, ул. „\makebox[10cm]{\dotfill}” № \makebox[0.7cm]{\dotfill}.}

\heading{Цели}
\article{Основни цели на СДРУЖЕНИЕТО са:}
\begin{itemize}
\item Да популяризира ходенето като едно от най-важните условия за {\bf изграждане на
  здраве};
\item Да способства ходенето пеша да се утвърди като личен и обществен избор, който
  допринася за {\bf чист въздух и осъзнатост};
\item Да помогне за изграждането на {\bf навици у нашите деца} за активен начин на
  живот;
\item Да припомни, че разходката е начин за споделяне на време и {\bf създаване на
  общност}.
\end{itemize}

\heading{Средства за постигане на целите}
\article{Целите на СДРУЖЕНИЕТО ще се осъществяват със следните средства:}
\begin{itemize}
\item Личен пример: Всеки член на СДРУЖЕНИЕТО носи отговорност както към себе си, така и
  към обкръжението си;
\item Организиране, провеждане и участие в дискусии, семинари, лекции и курсове на теми,
  свързани с целите на СДРУЖЕНИЕТО;
\item Организиране на разходки и спортни събития като начин за приобщаване и
  популяризиране на принципите на здравословен начин на живот;
\item Управление на проекти и кандидатстване за тяхното финансиране чрез международни
  или български програми, съобразени с целите на СДРУЖЕНИЕТО;
\item Участие в инициативи и проекти на други неправителствени организации, граждански
  институции, университети и болници;
\item Разработка на механизми за обмен на информация между различни институции и
  граждани от всички социални групи на обществото;
\item Изграждане на устойчиви механизми за диалог с държавната администрация и
  насърчаване на пряко и конструктивно участие на гражданите в този диалог;
\item Участие при планиране на изграждане на пешеходни пътни участъци и достъпен градски
  транспорт;
\item Иницииране и подпомагане на активни граждански дейности;
\item Популяризиране на целите и идеите на СДРУЖЕНИЕТО чрез брошури, електронни
  средства, публикации и кампании в медиите;
\item Създаване на прозрачна среда за комуникация и вземане на вътрешни за СДРУЖЕНИЕТО
  решения -- основополагащ елемент на конструктивно общуване с други сходни по идеи
  организации;
\item Използване на средства от членски внос и дарения, от стопанска дейност, доброволен
  труд и други дейности, разрешени от закона;
\end{itemize}

\heading{Предмет на дейност}

\article{}
\begin{enumerate}
\item Предметът на дейност на СДРУЖЕНИЕТО е осъществяване на дейности в обществена
  полза, насочени към популяризиране на ходенето и активния начин на живот като средство
  за изграждане и поддържане на здраве, опазване на околната среда и изграждане на
  активна гражданска общност.

\item СДРУЖЕНИЕТО организира разходки, спортни и образователни събития, обучения,
  семинари, курсове, кампании и инициативи. Кандидатства по, участва в, и управлява
  проекти, както и сътрудничи на други организации и институции.

\item СДРУЖЕНИЕТО участва в процеси по планиране и обсъждане на изграждането на
  пешеходна инфраструктура и достъпен градски транспорт и също така разпространява
  информация и материали, свързани с целите си.

\item СДРУЖЕНИЕТО може да създава или да участва в създаването у нас и в чужбина на
  самостоятелни и съвместни специализирани звена, представителства, клонове,
  изследователски центрове, клубове, фондове и сдружения с нестопанска цел, по реда,
  определен от закона.

\item СДРУЖЕНИЕТО може да извършва и допълнителна стопанска дейност, свързана с предмета
  и целите си, при условията на закона.
\end{enumerate}

\heading{Стопанска дейност}
\article{СДРУЖЕНИЕТО може да извършва допълнителна стопанска дейност, свързана с
  предмета и целите му, при условията на закона. Като стопанска дейност се включват:}
\begin{enumerate}
  \item Организиране на платени обучения, семинари, курсове и други събития, свързани с
    целите на СДРУЖЕНИЕТО.
  \item Рекламна, издателска и информационна дейност, включително производство и
    продажба на материали, свързани с целите на сдружението.
  \item Консултантски и посреднически услуги, включително предоставяне на платени
    обучения и услуги на трети лица.
  \item Управление на проекти на други организации, свързани с предмета и целите на
    СДРУЖЕНИЕТО.
  \item Отдаване под наем на движимо и недвижимо имущество.
  \item Предоставяне на услуги, свързани с информационни и комуникационни технологии,
    софтуерни решения и офис услуги.
  \item Организиране на платени културни, спортни и развлекателни събития.
  \item СДРУЖЕНИЕТО може да участва в създаването и управлението на търговски дружества
    и други стопански субекти, за подпомагане на своите цели, при условията на закона.
  \item Както и всяка друга стопанска дейност, свързана с визираните по-горе цели на
    СДРУЖЕНИЕТО.
\end{enumerate}
Приходите от стопанската дейност се използват единствено за постигане на целите на
СДРУЖЕНИЕТО.

\heading{Източници на средства}
\article{Средствата за осъществяване на дейността на СДРУЖЕНИЕТО се набавят от различни
  източници, включително:}
\begin{enumerate}
\item Членски внос на членовете на сдружението;
\item Дарения, завещания, спонсорства и други доброволни вноски от физически и
  юридически лица;
\item Приходи от проекти и програми, финансирани от национални или международни донори;
\item Приходи от допълнителна стопанска дейност, извършвана в рамките на закона;
\item Други източници, включително финансови инструменти, допустими по закон.
\end{enumerate}
Всички средства се използват изключително за постигане на целите на СДРУЖЕНИЕТО.

\heading{Определяне на дейността}
\article{СДРУЖЕНИЕТО ще осъществява дейност в ОБЩЕСТВЕНА ПОЛЗА по смисъла на {\bf Чл.
    2}, ал.~1 от ЗЮЛНЦ.}

\heading{Срок}
\article{СДРУЖЕНИЕТО не е ограничено със срок.}

\section{Членство}

\heading{Членство}
\article{Членове на СДРУЖЕНИЕТО могат да бъдат дееспособни физически лица и/или
  юридически лица, които заявят желание да спомогнат за постигане на целите и развитието
  на дейността на СДРУЖЕНИЕТО.}

\heading{Придобиване на членство}
\article{}
\begin{enumerate}
\item Членството в СДРУЖЕНИЕТО е доброволно.
\item Кандидатът подава заявление до Управителния съвет на СДРУЖЕНИЕТО, в коeто
  декларира, че е запознат и приема разпоредбите на настоящия Устав. Кандидатите --
  юридически лица представят заявление и преписи от документите си за регистрация и от
  решението на управителните си органи за членство в СДРУЖЕНИЕТО.
\item Управителният Съвет разглежда заявлението за членство на кандидата в срок от 7
  работни дни от датата на постъпването ѝ и с обикновено мнозинство (50\% плюс 1 гласа)
  я приема или отхвърля. Отказът на Управителния съвет може да се обжалва пред Общото
  събрание на СДРУЖЕНИЕТО. Членството в СДРУЖЕНИЕТО се придобива от датата на решението
  на Управителния съвет или от решението на Общото събрание.
\end{enumerate}

\heading{Права и задължения на членовете}
\article{Членовете на СДРУЖЕНИЕТО имат следните права:}
\begin{itemize}
\item да участват в управлението на СДРУЖЕНИЕТО, като избират Управителен съвет от 3
  лица;
\item да бъдат информирани за неговата дейност;
\item да ползват имуществото на СДРУЖЕНИЕТО за постигане целите на сдружението;
\item да ползват резултатите от дейността на СДРУЖЕНИЕТО съгласно разпоредбите на този
  Устав.
\end{itemize}

\article{Членовете на СДРУЖЕНИЕТО са длъжни:}
\begin{itemize}
\item да спазват разпоредбите на този Устав и да изпълняват решенията на ръководните
  органи на СДРУЖЕНИЕТО;
\item да участват в дейността на СДРУЖЕНИЕТО и да работят за осъществяване на целите му;
\item да издигат авторитета на СДРУЖЕНИЕТО;
\item да не извършват действия или бездействия, които противоречат на целите му, или да
  го злепоставят;
\item да внасят в срок предвидените в настоящия Устав членски вноски.
\end{itemize}

\article{Членските права и задължения, с изключение на имуществените, са непрехвърлими и
  не преминават върху други лица при смърт, съответно при прекратяване.}

\vspace{0.4cm}
\article{Членовете на СДРУЖЕНИЕТО имат право да овластят трето лице да упражнява техните
  права по конкретен повод и да изпълнява задълженията им, което се извършва писмено и
  произвежда действие само след писмено уведомяване на Управителния съвет. В този случай
  те носят отговорност за неизпълнение или изпълнение на техните задължения от страна на
  овластеното лице.}

\vspace{0.4cm}
\article{За задълженията на СДРУЖЕНИЕТО неговите членове носят отговорност само с
  предвидените в настоящия Устав членски вноски и кредиторите нямат право да предявяват
  правата си към личното им имущество над този размер.}

\heading{Прекратяване на членството}
\article{Членството в СДРУЖЕНИЕТО се прекратява:}
\begin{itemize}
\item с едностранно писмено волеизявление, отправено до Управителния съвет на
  СДРУЖЕНИЕТО;
\item със смърт или поставяне под запрещение, респ. прекратяване на юридическата личност
  на член на СДРУЖЕНИЕТО;
\item с изключване;
\item с прекратяване на СДРУЖЕНИЕТО.
\end{itemize}

\article{}
\begin{enumerate}
\item Член на СДРУЖЕНИЕТО може да бъде изключен с решение на Управителния съвет по
  негова инициатива или след постъпване на писмено предложение от 1/3 от членовете на
  СДРУЖЕНИЕТО, когато:
  \begin{itemize}
  \item нарушава предвидените в Устава задължения;
  \item извършва други действия, които правят по-нататъшното му членство в СДРУЖЕНИЕТО
    несъвместимо.
  \end{itemize}
\item При по-маловажни случаи на нарушения по {\bf Чл. 12} Управителният съвет определя
  с решение срок за преустановяване на нарушението и за отстраняването на неговите
  последици.
\end{enumerate}

\article{При прекратяване на членството СДРУЖЕНИЕТО не дължи връщане на направените
  членски вноски.}

\section{Управление}

\heading{Органи на управление}
\article{Върховен орган на СДРУЖЕНИЕТО е Общото събрание, негов орган е Управителният
  съвет.}

\heading{Общо събрание}
\article{}
\begin{enumerate}
\item Общото събрание се състои от всички членове на СДРУЖЕНИЕТО.
\item Юридическите лица участват в заседанието на Общото събрание чрез лицата, които ги
  представляват с изрично писмено нотариално заверено пълномощно. Пълномощниците нямат
  право да представляват повече от един член едновременно и да преупълномощават трети
  лица с правата си.
\end{enumerate}

\heading{Компетентност на общото събрание}
\article{Общото събрание:}
\begin{itemize}
\item Изменя и допълва Устава.
\item Избира и освобождава членовете  на Управителния съвет.
\item Взема решение за участие в други организации.
\item Взема решение за преобразуване или прекратяване на СДРУЖЕНИЕТО.
\item Приема основните насоки и програми за дейността на СДРУЖЕНИЕТО.
\item Приема бюджета на СДРУЖЕНИЕТО.
\item Произнася се с решение по жалби на изключени и кандидат-членове.
\item Приема отчета за дейността на Управителния съвет.
\item Отменя решенията на другите органи на СДРУЖЕНИЕТО, които противоречат на закона,
  Устава или други вътрешни актове, регламентиращи дейността на СДРУЖЕНИЕТО.
\item Взема решения по всички други въпроси, поставени в неговата компетентност от
  закона или настоящия Устав.
\end{itemize}

\heading{Свикване}
\article{Общото събрание се свиква на редовно заседание от Управителния съвет на
  СДРУЖЕНИЕТО, веднъж годишно. Заседанието може да се проведе в подходящо помещение в
  населеното място по седалище на СДРУЖЕНИЕТО или онлайн (дистанционно).}

\vspace{0.4cm}
\article{Извънредно заседание на Общо събрание се свиква по инициатива на Управителния
  съвет. Една трета от членовете на СДРУЖЕНИЕТО имат право да поискат Управителния съвет
  да свика Общо събрание по конкретен въпрос и ако той не отправи покана в едномесечен
  срок от датата на поискването, събранието се свиква от съда по седалището на
  СДРУЖЕНИЕТО.}

\heading{Кворум}
\article{}
\begin{enumerate}
\item Общото събрание е законно, ако присъстват повече от половината от всички членове.
\item Кворумът се установява от председателстващия събранието и секретаря по списък, в
  който се отразяват имената на присъстващите членове и техните представители, подписва
  се от тях, заверява се от председателя и секретаря на заседанието и се прилага към
  протокола за него.
\end{enumerate}

\heading{Гласуване}
\article{При гласуване всеки член на СДРУЖЕНИЕТО има право само на един глас.}

\heading{Конфликт на интереси}
\article{Член на СДРУЖЕНИЕТО няма право на глас при решаване на въпроси отнасящи се до:}
\begin{itemize}
\item него, съпруг/а или роднина по права линия - без ограничения, по съребрена линия -
  до четвърта степен, или по сватовство - до втора степен включително;
\item юридически лица, в които той е управител или може да наложи или възпрепятства
  вземането на решения.
\end{itemize}

\heading{Решения}
\article{Решенията на Общото събрание се вземат с обикновено мнозинство (50\% плюс 1
  гласа) от присъстващите. Решенията по {\bf Чл. 21}, се вземат с квалифицирано
  мнозинство 2/3 от присъстващите.}

\vspace{0.4cm}
\article{Общото събрание не може да взема решения, които не са включени в предварително
  обявения дневен ред на събранието.}

\heading{Протокол}
\article{}
\begin{enumerate}
\item За всяко заседание на Общото събрание се води протокол, който се заверява от
  председателстващия събранието и лицето, изготвило протокола, които отговарят за
  верността на съдържанието му.
\item Всеки член, присъствал на Общото събрание има право да следи за точното отразяване
  на заседанието и взетите решения в протокола.
\end{enumerate}

\heading{Управителен съвет}
\article{Управителният съвет се състои от 3 лица -- членове на СДРУЖЕНИЕТО.}

\heading{Мандат}
\article{Управителният съвет се избира за срок от 3 години, като неговите членове могат
  да бъдат преизбирани неограничено.}

\heading{Правомощия}
\article{Управителният съвет:}
\begin{itemize}
\item Представлява СДРУЖЕНИЕТО и определя обема на представителната власт, на
  председателя и секретаря.
\item Избира из между членовете си председател и секретар.
\item Осигурява изпълнението на решенията на Общото събрание.
\item Взема решение за откриване и закриване на клонове и назначава управител/и.
\item Подготвя и внася в Общото събрание проект за бюджет.
\item Подготвя и внася в Общото събрание отчет за дейността на СДРУЖЕНИЕТО.
\item Определя реда и организира извършването на дейността на СДРУЖЕНИЕТО и носи
  отговорност за това.
\item Определя адреса на СДРУЖЕНИЕТО.
\item Приема правила за работата си и пряко ръководи стопанската дейност на СДРУЖЕНИЕТО.
\item Взема решения по всички въпроси, освен тези които са от компетентността на Общото
  събрание на СДРУЖЕНИЕТО.
\item Приема и изключва членове на СДРУЖЕНИЕТО.
\end{itemize}

\heading{Заседания}
\article{Редовните заседанията на Управителния съвет се свикват от председателя по
  негова инициатива веднъж месечно, както и по искане на всеки един от неговите членове.
  Ако председателят не свика заседание в срок от 7 работни дни от поискването, такова
  може да се свика от всеки един от членовете на Управителния съвет.}

\vspace{0.4cm}
\article{}
\begin{enumerate}
\item Заседанието е законно, ако на него присъстват повече от половината от членовете на
  Управителния съвет.
\item Законно решение може да бъде взето и без да се провежда заседание, ако протоколът
  за това бъде подписан без забележки и възражения за това от всички членове на
  Управителния съвет.
\item Относно протоколите за заседанията на Управителния съвет се прилагат съответно
  разпоредбите на {\bf Чл. 31}.
\end{enumerate}

\article{Заседанията се ръководят от председателя на Управителния съвет, а в негово
  отсъствие от посочен от него член.}

\heading{Решения}
\article{Управителният съвет взема решенията си с мнозинство от присъстващите членове.}

\heading{Контрол}
\article{Контрол се осъществява от Общото събрание на СДРУЖЕНИЕТО.}

\heading{Отговорност на членовете на управителния съвет}
\article{Членовете на Управителния съвет носят солидарна отговорност за действията си,
  увреждащи имуществото и интересите на СДРУЖЕНИЕТО.}

\heading{Председател на управителния съвет}
\article{Управителният cъвет избира от своя състав председател, като с решението си
  определя и неговите функции.}

\heading{Клонове}
\article{С решение на Управителния съвет на СДРУЖЕНИЕТО могат да се откриват и закриват
  клонове извън населеното място, в което се намира седалището на СДРУЖЕНИЕТО.}

\vspace{0.4cm}
\article{Клоновете не са юридически лица, ръководят се от управител, назначен от
  Управителния съвет и извършват дейностите, определени с решението на Управителния
  съвет. Със същото решение се определят и ограниченията в правомощията им и
  представителната власт на управителя на клона.}

\vspace{0.4cm}
\article{Клоновете водят книги за дейността си и веднъж на месец управителят на клона
  представя пред Управителният съвет на СДРУЖЕНИЕТО отчет за дейността на клона и
  разходваните средства.}

\vspace{0.4cm}
\article{Управителният съвет на СДРУЖЕНИЕТО заявява пред съда, района в който се намира
  седалището на клона, наименованието на седалището и адреса на СДРУЖЕНИЕТО, седалището
  и адреса на клона, неговия управител и ограниченията на неговите правомощия и
  представителна власт. На заявление подлежат и промените в горепосочените
  обстоятелства. Заявлението се подава в срок от 7 работни дни от датата на решение на
  Управителния съвет.}

\heading{Членски внос}
\article{Членският внос се определя в размер на 25 евро за физически лица и половин
  минимална работна заплата за юридически лица, а в годината на учредяване - 1/3 от
  дължимия годишен членски внос.}

\section{Прекратяване и ликвидация}

\heading{Прекратяване}
\article{}
\begin{itemize}
\item по решение на Общото събрание;
\item по решение на Окръжния съд по седалището в случаите по {\bf Чл.~13}, ал.~1, т.~3
  от ЗЮЛНЦ.
\end{itemize}

\section{Заключителни разпоредби}
\begin{itemize}
\item Този устав е приет на учредително събрание за създаване на сдружение с нестопанска
  цел „ПЕШАЧКО”, проведено на \makebox[5cm]{\dotfill}. г. в гр. София, \\ул.
  „\makebox[9cm]{\dotfill}” № \makebox[0.7cm]{\dotfill}.
\item Списъкът на членовете на СДРУЖЕНИЕТО, подписали устава да се счита за неразделна
  част от този Устав.
\end{itemize}

\vspace{1cm}
{\bf\large СПИСЪК НА ЧЛЕНОВЕТЕ НА НПО „ПЕШАЧКО”}: \\[0.5cm]

\begin{itemize}
\item \makebox[10cm]{\dotfill}, ЕГН: \makebox[4.8cm]{\dotfill} \\[0.5cm]
\item \makebox[10cm]{\dotfill}, ЕГН: \makebox[4.8cm]{\dotfill} \\[0.5cm]
\item \makebox[10cm]{\dotfill}, ЕГН: \makebox[4.8cm]{\dotfill} \\[0.5cm]
\item \makebox[10cm]{\dotfill}, ЕГН: \makebox[4.8cm]{\dotfill} \\[0.5cm]
\item \makebox[10cm]{\dotfill}, ЕГН: \makebox[4.8cm]{\dotfill} \\[0.5cm]
\item \makebox[10cm]{\dotfill}, ЕГН: \makebox[4.8cm]{\dotfill} \\[0.5cm]
\item \makebox[10cm]{\dotfill}, ЕГН: \makebox[4.8cm]{\dotfill}
\end{itemize}

\end{document}
